\subsection{\texorpdfstring{SYSTEM OF TYPE \(F_{4}\)}{SYSTEM OF TYPE F\_4}}

\begin{enumerate}
\def\labelenumi{(\Roman{enumi})}
\tightlist
\item
  Consider the group \(\mathbf{L}_2\) (no. 4) in \(\mathbf{R}^4\) . Let \(\mathbf{R}\) be the set of \(\alpha \in \mathbf{L}_2\) such that \((\alpha |\alpha) = 1\) or \((\alpha |\alpha) = 2\) ; it contains the vectors
\end{enumerate}

\[
\pm \varepsilon_ {i}, \quad \pm \varepsilon_ {i} \pm \varepsilon_ {j} (i <   j), \quad \frac {1}{2} (\pm \varepsilon_ {1} \pm \varepsilon_ {2} \pm \varepsilon_ {3} \pm \varepsilon_ {4}).
\]

Conversely, if \(\alpha\in R\) , the coordinates of \(\alpha\) can only take the values \(0,\pm\frac{1}{2},\pm1\) (since \((\frac{3}{2})^{2}>2\) ); either these coordinates are all integers, giving the vectors \(\pm\varepsilon_{i},\pm\varepsilon_{i}\pm\varepsilon_{j}\) , or they are all equal to \(\pm\frac{1}{2}\) , giving the vectors \(\frac{1}{2}(\pm\varepsilon_{1}\pm\varepsilon_{2}\pm\varepsilon_{3}\pm\varepsilon_{4})\) .

We show that, for \(\alpha, \beta \in \mathbb{R}\) , we have \(2(\alpha|\beta) / (\alpha|\alpha) \in \mathbf{Z}\) . If \(\alpha = \pm \varepsilon_i\) or \(\alpha = \frac{1}{2} (\pm \varepsilon_1 \pm \varepsilon_2 \pm \varepsilon_3 \pm \varepsilon_4)\) , then \((\alpha|\alpha) = 1\) and we have seen in no. 4 that \((\alpha|\beta) \in \frac{1}{2}\mathbf{Z}\) since \(\alpha, \beta \in \mathbf{L}_2\) . If \(\alpha = \pm \varepsilon_i \pm \varepsilon_j\) , then \((\alpha|\alpha) = 2\) and we have seen in no. 4 that \((\alpha|\beta) \in \mathbf{Z}\) since \(\alpha \in \mathbf{L}_1\) and \(\beta \in \mathbf{L}_2\) . Hence, R is a reduced root system in \(\mathbf{R}^4\) (no. 4). The number of roots is \(n = 8 + \binom{4}{2} 4 + 2^4 = 48\) .

\begin{enumerate}
\def\labelenumi{(\Roman{enumi})}
\setcounter{enumi}{1}
\tightlist
\item
  Give \(\mathbf{R}^4\) the lexicographic order defined by the basis \((\varepsilon_1, \varepsilon_2, \varepsilon_3, \varepsilon_4)\) (§ 1, no. 7). In particular, we have \(\varepsilon_1 > \varepsilon_2 > \varepsilon_3 > \varepsilon_4\) . The positive roots are
\end{enumerate}

\[
\varepsilon_ {i}, \quad \varepsilon_ {i} \pm \varepsilon_ {j} (i <   j), \quad \frac {1}{2} (\varepsilon_ {1} \pm \varepsilon_ {2} \pm \varepsilon_ {3} \pm \varepsilon_ {4}).
\]

The smallest root is \(\alpha_{3} = \varepsilon_{4}\) . Among the roots belonging to \(\mathbf{R}\varepsilon_{3} + \mathbf{R}\varepsilon_{4}\) but not to \(\mathbf{R}\varepsilon_{4}\) , the smallest is \(\alpha_{2} = \varepsilon_{3} - \varepsilon_{4}\) . Among the positive roots belonging to \(\mathbf{R}\varepsilon_{2} + \mathbf{R}\varepsilon_{3} + \mathbf{R}\varepsilon_{4}\) but not to \(\mathbf{R}\varepsilon_{3} + \mathbf{R}\varepsilon_{4}\) , the smallest is \(\alpha_{1} = \varepsilon_{2} - \varepsilon_{3}\) . Among the positive roots not belonging to \(\mathbf{R}\varepsilon_{2} + \mathbf{R}\varepsilon_{3} + \mathbf{R}\varepsilon_{4}\) , the smallest is \(\alpha_{4} = \frac{1}{2} (\varepsilon_{1} - \varepsilon_{2} - \varepsilon_{3} - \varepsilon_{4})\) . None of the \(\alpha_{i}\) is a sum of 2 positive roots. Hence, \((\alpha_{1},\alpha_{2},\alpha_{3},\alpha_{4})\) is a basis of R (§ 1, no. 6, Cor. 1 to Prop. 19). We have \(\| \alpha_1 \|^2 = \| \alpha_2 \|^2 = 2\) , \(\| \alpha_3 \|^2 = \| \alpha_4 \|^2 = 1\) , \((\alpha_1|\alpha_2) = (\alpha_2|\alpha_3) = -1\) , \((\alpha_3|\alpha_4) = -\frac{1}{2}\) , \((\alpha_1|\alpha_3) = (\alpha_1|\alpha_4) = (\alpha_2|\alpha_4) = 0\) . We see that the Dynkin graph of R is of type \(F_4\) , and hence is irreducible.

\begin{enumerate}
\def\labelenumi{(\Roman{enumi})}
\setcounter{enumi}{2}
\item
  We have \(h = \frac{n}{l} = 12\) .
\item
  Let \(\tilde{\alpha} = \varepsilon_1 + \varepsilon_2 = 2\alpha_1 + 3\alpha_2 + 4\alpha_3 + 2\alpha_4\) . The sum of the coordinates of \(\tilde{\alpha}\) with respect to \((\alpha_i)\) is \(11 = h - 1\) , so \(\tilde{\alpha}\) is the highest root. We have \((\tilde{\alpha}|\alpha_1) = 1\) , \((\tilde{\alpha}|\alpha_2) = (\tilde{\alpha}|\alpha_3) = (\tilde{\alpha}|\alpha_4) = 0\) .
\end{enumerate}

The completed Dynkin graph is

\[
\alpha_ {1} \quad \alpha_ {2} \quad \alpha_ {3} \quad \alpha_ {4}
\]

\begin{enumerate}
\def\labelenumi{(\Alph{enumi})}
\setcounter{enumi}{21}
\tightlist
\item
  The formula \(\alpha^{\prime} = \frac{2\alpha}{(\alpha|\alpha)}\) gives for \(R^{\prime}\) the set of vectors \(\pm 2\varepsilon_{i}, \pm \varepsilon_{i} \pm \varepsilon_{j}, \pm \varepsilon_{1} \pm \varepsilon_{2} \pm \varepsilon_{3} \pm \varepsilon_{4}\) . The Dynkin graph of \(R^{\prime}\) is obtained from that of \(R\) by the procedure explained in no. 2. and we see that \(R^{\prime}\) is of type \(F_{4}\) .
\end{enumerate}

The roots not orthogonal to \(\beta = \varepsilon_1\) are \(\pm \varepsilon_1, \pm \varepsilon_1 \pm \varepsilon_j\) ( \(j \geqslant 2\) ), and \(\frac{1}{2} (\pm \varepsilon_1 \pm \varepsilon_2 \pm \varepsilon_3 \pm \varepsilon_4)\) ; the number \(n(\alpha | \beta) = 2(\alpha | \beta)\) is equal to \(\pm 2\) for the first 14 of these roots and to \(\pm 1\) for the last 16; thus, for \(\Phi_R\) , the square of the length of \(\beta\) is \(4(14.4 + 16.1)^{-1} = \frac{1}{18}\) ; hence

\[
\Phi_ {R} (x, y) = \frac {(x | y)}{18}.
\]

We now apply formula (18) of § 1, no. 12, with \(x = y = \beta\) ; this gives

\[
2 + 12. \frac {1}{4} + 16. \frac {1 / 4}{1} = \gamma (\mathrm{R}). \frac {1}{18}
\]

so

\[
\gamma (\mathrm{R}) = 2. 3 ^ {4}.
\]

\begin{enumerate}
\def\labelenumi{(\Roman{enumi})}
\setcounter{enumi}{5}
\tightlist
\item
  Calculating the fundamental weights gives
\end{enumerate}

\[
\omega_ {1} = \varepsilon_ {1} + \varepsilon_ {2} = 2 \alpha_ {1} + 3 \alpha_ {2} + 4 \alpha_ {3} + 2 \alpha_ {4} = \tilde {\alpha}
\]

\[
\omega_ {2} = 2 \varepsilon_ {1} + \varepsilon_ {2} + \varepsilon_ {3} = 3 \alpha_ {1} + 6 \alpha_ {2} + 8 \alpha_ {3} + 4 \alpha_ {4}
\]

\[
\omega_ {3} = \frac {1}{2} (3 \varepsilon_ {1} + \varepsilon_ {2} + \varepsilon_ {3} + \varepsilon_ {4}) = 2 \alpha_ {1} + 4 \alpha_ {2} + 6 \alpha_ {3} + 3 \alpha_ {4}
\]

\[
\omega_ {4} = \varepsilon_ {1} = \alpha_ {1} + 2 \alpha_ {2} + 3 \alpha_ {3} + 2 \alpha_ {4}.
\]

\begin{enumerate}
\def\labelenumi{(\Roman{enumi})}
\setcounter{enumi}{6}
\tightlist
\item
  The sum of the positive roots is
\end{enumerate}

\[
2 \rho = 11 \varepsilon_ {1} + 5 \varepsilon_ {2} + 3 \varepsilon_ {3} + \varepsilon_ {4} = 16 \alpha_ {1} + 30 \alpha_ {2} + 42 \alpha_ {3} + 22 \alpha_ {4}.
\]

\begin{enumerate}
\def\labelenumi{(\Roman{enumi})}
\setcounter{enumi}{7}
\item
  We have \(\mathrm{Q}(\mathrm{R}) = \mathrm{L}_2\) (no. 4), and \(\mathrm{P}(\mathrm{R}) = \mathrm{Q}(\mathrm{R})\) by (VI). Hence, the connection index is 1.
\item
  The family of exponents has 4 terms, and since \(h = 12\) , it must contain the integers 1, 5, 7, 11, coprime to 12 (§ 1, no. 11, Prop. 30); consequently, these are all the exponents of \(W(R)\) .
\item
  (X) and (XI) The only automorphism of the Dynkin graph is the identity, so \(A(R) = W(R)\) and \(w_{0} = -1\) . Let \(R'\) be the set of longest elements of R, that is, the \(\pm\varepsilon_{i} \pm \varepsilon_{j}\) : \(R'\) is the root system of type \(D_{4}\) constructed in no. 8. Clearly, every element of \(A(R)\) is an element of \(A(R')\) . Conversely, an element of \(A(R')\) leaves \(L_{1}\) stable (since it is generated by \(R'\) ), hence also its associated \(L_{2}\) , and hence also R. So \(W(R) = A(R) = A(R')\) . By no. 8, \(W(R)\) is the semi-direct product of \(S_{3}\) by \(W(R')\) , \(W(R')\) itself being the semi-direct product of \(S_{4}\) by \((\mathbf{Z}/2\mathbf{Z})^{3}\) . The order of \(W(R)\) is \(3!4!2^{3} = 2^{7}.3^{2}\) .
\end{enumerate}

